\documentclass[11pt]{article}
\usepackage{setspace}
\usepackage{geometry}
\geometry{margin=1in}
\setstretch{1.15}

\begin{document}

% --- Applicant Information ---
\noindent
\textbf{Zhongli Wang} \\
Phone: +49 160 9541 6015 \\
Email: nkwangzhongli@gmail.com \\
Website: https://sites.google.com/view/zhongli-wang \\
Bielefeld, Germany \\

\vspace{0.5cm}

% --- Recipient Information ---
\noindent
\textbf{Recruitment Committee} \\
Faculty of Business, Economics and Social Sciences \\
Department of Social Sciences \\
Journalism and Communication Studies (Digital Journalism) \\
Universit\"at Hamburg \\
Hamburg, Germany \\

\vspace{0.5cm}

% --- Subject Line ---
\noindent
\textbf{Application for Research Associate (Postdoc) --- Project ``PSS4DEMOCRACY''} \\

\vspace{0.3cm}

\noindent Dear Prof. Lischka and Members of the Selection Committee,\\[0.2cm]

I am writing to apply for the Research Associate (Postdoc) position in the project \textit{``Which Public Sphere System for Democracy? -- PSS4DEMOCRACY''} at Universit\"at Hamburg. I am currently completing my Ph.D.\ in Economics at Bielefeld University and Universit\`a Cattolica del Sacro Cuore through the Marie Sk{\l}odowska--Curie EPOC programme. My work sits at the intersection of computational social science and applied modeling: I build reproducible, data-intensive pipelines to measure latent concepts from text and to study how algorithmic environments shape human behavior. The advertised position---translating normative public-sphere ideals into operational NLP metrics and experimentally testable algorithmic news feeds---is a strong match to my methodological background and research interests.

A central theme of my doctoral research is how information environments and social influence shape decisions under heterogeneity and bounded rationality. In my dissertation, \textit{``Strategic Behavior and Economic Dynamics in ICT Markets: An Agent-Based Framework,''} I develop computational models in which heterogeneous agents update beliefs, react to information frictions, and generate aggregate outcomes through interaction and feedback. While my substantive application is digital markets, the core methodological problem is closely aligned with PSS4DEMOCRACY: formalizing abstract objectives (e.g., diversity, exposure, salience, emotionalization) into measurable constructs and then evaluating system-level consequences under alternative algorithmic rules.

Beyond modeling, I have extensive experience building end-to-end, reproducible computational pipelines in Python. In a coauthored project on \textit{agentic workflows for economic research}, I developed a framework that uses large-language-model agents to automate and document multi-stage research tasks (data collection, coding, text processing, and analysis) with strong emphasis on traceability, version control, and reproducibility. This experience translates directly to the computational core of PSS4DEMOCRACY: implementing a reliable NLP-to-ranking pipeline, maintaining experiment-ready data products, and producing transparent metric specifications and validation reports that can be audited by an interdisciplinary team.

I am particularly excited about the project\textquotesingle s two main components. First, on \textit{NLP-based news content metrics}, I can contribute to implementing and validating multilingual pipelines for entity-centric and viewpoint-sensitive measures. Concretely, I am comfortable operationalizing concepts through modular components such as entity linking and disambiguation, topic modeling, stance- and sentiment/emotion-oriented classifiers, and document-level aggregation into interpretable indices (e.g., party visibility and issue attention shares, distributional measures of viewpoint diversity, or calibrated emotionalization intensity). I also have experience designing validation protocols that compare automated measures to human-coded labels via reliability and error diagnostics, and iterating models based on where measurement fails.

Second, on \textit{news recommender system development}, I am motivated by the project\textquotesingle s emphasis on normative multi-objective ranking. My background in simulation and optimization has trained me to think in terms of objective trade-offs, constraints, and counterfactual evaluation. I can contribute to implementing multi-objective scoring and ranking approaches (e.g., scalarization baselines, Pareto-front ranking, and constrained optimization variants), and to setting up robust offline evaluation with standard recommendation metrics alongside public-sphere objective metrics. While my prior work is not in production recommender systems per se, I have strong experience in building scalable experiments, running ablations, and documenting algorithmic choices in a way that supports credible scientific claims and replication.

I also value the collaborative and coordination aspects of the role. In my doctoral programme, I have worked in international teams and coordinated research workflows that involve shared codebases, standardized data schemas, and iterative evaluation. I would be glad to take responsibility for maintaining the computational infrastructure (data processing, model training, evaluation dashboards, and reproducibility protocols), and to support project coordination by translating between normative theory, measurement design, and implementable engineering choices.

Universit\"at Hamburg\textquotesingle s Digital Journalism group offers an ideal setting for this work: a strong normative and empirical research environment where algorithmic systems are studied not only as engineering artifacts but as democratic institutions with measurable, testable implications. I would welcome the opportunity to contribute my computational expertise to PSS4DEMOCRACY and to coauthor publications that advance both methodological rigor and the project\textquotesingle s core democratic aims.

Thank you very much for considering my application. I would be delighted to discuss how my background in computational social science, NLP-oriented measurement, and reproducible pipeline development can support the goals of PSS4DEMOCRACY.

\vspace{6mm}

\noindent Sincerely,\\[12mm]
\noindent Zhongli Wang

\end{document}